\begin{pastexam}[label={px:df3-101-ds-04}]{地方三等資訊處理 101 年 第四題}
（申論；20 分）

迴文（palindrome）乃是一個字串不論從左到右或從右到左看結果一模一樣，例如“油麻地遍地麻油”、“人人為我、
我為人人”、“Fall leaves as soon as leaves fall”等。
\begin{itemize}[leftmargin=2.2em]
\item[㈠] 判斷是否為迴文，該選用何種資料結構或抽象資料型態來儲存字串，為什麼？（5 分）
\item[㈡] 請寫出演算法以測試輸入的字串（字串的長度不超過 100 個字）是否為迴文。（10 分）
\item[㈢] 請分析該演算法的時間複雜度（time complexity）。（5 分）
\end{itemize}
\end{pastexam}

\begin{pastsoln}
答案：㈠ 用堆疊（字串本身存在一維陣列裡）；㈡ 前半段推入堆疊，再一邊彈出一邊和後半段比較；
㈢ $O(n)$。 本題歸第~6~章。

先釐清「一模一樣」比的單位。 前兩個中文例子逐字倒過來確實一樣（「、」也在正中間）。 第三個英文例子若逐個
字母倒過來是 “llaf sevael sa noos ...”，並不一樣；它是\textbf{以單字為單位}倒過來相同（Fall / leaves /
as / soon / as / leaves / fall），而且不分大小寫。 所以下面的演算法先把字串切成「單位」：每個中文字一個
單位，每個英文單字（轉成小寫）一個單位，空白和標點不算，再判斷單位序列是不是迴文。 只考慮逐字元的迴文時，
把每個字元當一個單位即可，演算法不變。

\textbf{㈠} 字串的單位依序存在一維陣列 $T[1..n]$（$n \le 100$，長度有上限，陣列大小固定 100 就夠）。
判斷時用\textbf{堆疊}：堆疊後進先出，把前半段依序推入，彈出時的順序正好是前半段的反序，拿來和後半段由左
往右逐一比較，相同就是迴文。 這和資結書第~6.1~節「用堆疊把序列反過來」是同一個用法。 （也可以只用陣列、
左右兩個索引往中間夾，或用雙向佇列從兩端各取一個比較；這些做法時間都一樣。）

\textbf{㈡} 設 $T[1..n]$ 已經是切好的單位序列。
\begin{pseudo}[unbreakable]{用堆疊判斷迴文}
\Function{Is-Palindrome}{$T, n$}
    \State \Call{Stack}{$S$}
    \State $h \gets \lfloor n/2 \rfloor$
    \For{$i \gets 1$ \To $h$}
        \State \Call{Push}{$S, T[i]$} \Comment{前半段推入}
    \EndFor
    \State $j \gets n - h + 1$ \Comment{$n$ 為奇數時跳過正中間 $T[h+1]$}
    \While{$j \le n$}
        \State \Call{Pop}{$S, x$}
        \If{$x \ne T[j]$}
            \State \Return \FALSE
        \EndIf
        \State $j \gets j + 1$
    \EndWhile
    \State \Return \TRUE
\EndFunction
\end{pseudo}
正確性：第 3--5 行後，堆疊由頂到底是 $T[h], T[h-1], \ldots, T[1]$。 第 7--11 行第 $t$ 次彈出的是
$T[h + 1 - t]$，比較的對象是 $T[n - h + t]$，兩個下標的和是 $n + 1$，也就是比較「從左數第 $k$ 個」和
「從右數第 $k$ 個」（$k = h + 1 - t$ 跑遍 $1..h$）。 迴文的定義正是這 $h$ 對都相等（$n$ 為奇數時正中間
那一個和自己比，不必檢查），所以全部相同回傳 \TRUE，有一對不同就回傳 \FALSE。

追蹤「人人為我、我為人人」（$n = 9$，$h = 4$）：推入 人、人、為、我；跳過正中間的「、」；從 $T[6]$ 起依序
比較 我／我、為／為、人／人、人／人，全部相同，是迴文。 以單字為單位的英文例子 $n = 7$、$h = 3$，推入
fall、leaves、as，跳過 soon，比較 as／as、leaves／leaves、fall／fall，也是迴文。

前處理（切單位）：從左到右掃描原字串一次，中文字直接成為一個單位，連續的英文字母合成一個單字並轉小寫，
空白和標點略過，也是一次掃描。

\textbf{㈢} 設字串長度（或單位數）是 $n$。 前處理掃描一次 $O(n)$；推入 $\lfloor n/2 \rfloor$ 次、彈出和比較
最多 $\lfloor n/2 \rfloor$ 次，每次 $O(1)$（資結書第~6.1~節），所以總時間 $\mathbf{O(n)}$。 最好情況是第一對
就不同，比較一次就結束，但前處理和推入仍要 $\Theta(n)$。 額外空間是堆疊的 $O(n)$；改用左右兩個索引往中間
夾的做法，額外空間只要 $O(1)$。 本題字串最多 100 個字，最多約 50 次推入、50 次比較。
\end{pastsoln}
